\documentclass[main.tex]{subfiles}

\begin{document}

The vast majority of recorded music available today is in stereophonic format. The \gls{sip} is a description of the use of stereo field in a musical recording or production. It is assessed based on the final recording without prior knowledge of how said recording was created. There are many established methods for creating a \gls{sip}, each useful for differing downstream tasks or goals, and each showcasing different spatial information present in the recording. This thesis will investigate \gls{sip} methods, gaining understanding of what spatial information can be extracted and what downstream tasks can be achieved by examining the \glspl{sip}.

The thesis is divided into 4 papers. Paper-A examines 4 existing methods for describing the use of amplitude panning in music production, highlighting the capabilities and limitations of each one. Paper-B presents a perceptual listening test of audio in which the spatial properties are varied. Paper-C is a much larger systematic review of all methods used to create \glspl{sip}, and their downstream uses. Paper-D details an experiment in which the \gls{sip} methods from previous papers are put to use to find music tracks that have similar uses of stereo field.

It is shown that while the established \gls{sip} methods can be useful for real-world tasks, they sometimes lack the ability to identify more subtle uses of stereo field. For example, they predominantly focus on the use of amplitude panning, but more modern spatial audio formats are typically rendered to a stereo binaural signal, in which both phase and amplitude differences between channels can be equally important. This thesis establishes a \gls{sip} baseline from which future research can be spring-boarded.

\end{document}